% ============================================================================
\section{Understanding Transformers and Large Language Models}
\label{sec:transformers-llms}
% ============================================================================

The previous sections developed four paradigms for replacing
floating-point arithmetic with bitwise operations, and showed how they
combine into a hybrid architecture for MNIST. Before we can honestly
evaluate whether any of this scales to large language models (LLMs),
we need to understand what LLMs actually \emph{are}---what makes them
work, and which of their components are essential versus replaceable.

This section builds that understanding from the ground up, using the
same intuition-first approach as the rest of this article.

% ----------------------------------------------------------------------------
\subsection{From Pixels to Words: Why Language Is Harder}
\label{sec:pixels-vs-words}
% ----------------------------------------------------------------------------

MNIST is a \emph{fixed-size classification} problem: 784 pixels in,
one of ten labels out. Every input has the same shape, and every pixel
has a clear spatial relationship to its neighbours.

Language is fundamentally different:

\begin{itemize}[leftmargin=2em]
  \item \textbf{Variable length.} A sentence can be 5~words or
    5{,}000~words. The network must handle both.
  \item \textbf{No fixed grid.} Words are not arranged in a 2D grid
    like pixels. Their meaning depends on \emph{order} and
    \emph{context}---``the dog bit the man'' and ``the man bit the
    dog'' use the same words but mean different things.
  \item \textbf{Long-range dependencies.} In ``The cat that the dog
    chased ran away,'' the verb ``ran'' relates to ``cat,'' not the
    nearby ``dog.'' The network must connect words across arbitrary
    distances.
  \item \textbf{Generation, not classification.} An LLM does not pick
    one of ten labels. It generates text one word at a time, choosing
    from a vocabulary of 30{,}000--100{,}000 possible next words at
    each step.
\end{itemize}

\begin{keyinsight}
The transition from MNIST to language is not just ``bigger data.'' It
requires fundamentally different architectural primitives: a way to
represent words as numbers, a mechanism to relate words to each other
regardless of distance, and a method for generating sequences one
element at a time.
\end{keyinsight}

% ----------------------------------------------------------------------------
\subsection{Tokenisation: Turning Text into Numbers}
\label{sec:tokenisation}
% ----------------------------------------------------------------------------

A neural network cannot process letters or words directly---it needs
numbers. The first step in any language model is \concept{tokenisation}:
splitting text into small pieces and assigning each piece an integer ID.

Modern LLMs do not split on whole words. Instead, they use
\concept{subword tokenisation} (algorithms like Byte-Pair Encoding or
SentencePiece), which breaks text into frequently occurring fragments:

\begin{figure}[htbp]
\centering
\begin{tikzpicture}[
  tok/.style={draw, rounded corners=2pt, minimum height=0.7cm,
    font=\small\ttfamily, inner sep=4pt, fill=#1!10, draw=#1!60},
  arr/.style={-{Stealth[length=4pt]}, thick, bitgray}
]
  % Original text
  \node[font=\small\sffamily] at (-1.5, 1.5) {Input text:};
  \node[font=\small\ttfamily] at (4, 1.5)
    {``Understanding transformers is fascinating''};

  % Tokens
  \node[font=\small\sffamily] at (-1.5, 0) {Tokens:};
  \node[tok=bitblue]   (t1) at (1.5, 0) {Under};
  \node[tok=bitgreen]  (t2) at (3.3, 0) {stand};
  \node[tok=bitorange] (t3) at (5.0, 0) {ing};
  \node[tok=bitpurple] (t4) at (7.0, 0) {\_transform};
  \node[tok=bitblue]   (t5) at (9.3, 0) {ers};
  \node[tok=bitgreen]  (t6) at (11.0, 0) {\_is};
  \node[tok=bitorange] (t7) at (13.0, 0) {\_fasc};
  \node[tok=bitpurple] (t8) at (14.8, 0) {inating};

  % Token IDs
  \node[font=\small\sffamily] at (-1.5, -1.2) {Token IDs:};
  \node[font=\small\ttfamily] at (1.5, -1.2) {8803};
  \node[font=\small\ttfamily] at (3.3, -1.2) {1514};
  \node[font=\small\ttfamily] at (5.0, -1.2) {278};
  \node[font=\small\ttfamily] at (7.0, -1.2) {13389};
  \node[font=\small\ttfamily] at (9.3, -1.2) {414};
  \node[font=\small\ttfamily] at (11.0, -1.2) {338};
  \node[font=\small\ttfamily] at (13.0, -1.2) {19239};
  \node[font=\small\ttfamily] at (14.8, -1.2) {15536};
\end{tikzpicture}
\caption{Subword tokenisation splits text into frequently occurring
  fragments, each mapped to an integer ID. Common words become single
  tokens; rare words are split into subwords. The underscore
  (\_) represents a space preceding the token.}
\label{fig:tokenisation}
\end{figure}

\begin{analogy}
Tokenisation is like Morse code for neural networks. Just as Morse
represents each letter as a sequence of dots and dashes, tokenisation
represents each text fragment as a number. The difference is that the
``alphabet'' is learned from data: common fragments (``the'', ``ing'',
``tion'') get their own codes, while rare fragments are spelled out
from smaller pieces. A typical LLM vocabulary has 30{,}000--100{,}000
tokens.
\end{analogy}

The key property is that \emph{any} text can be represented as a
sequence of token IDs. The model's input is always a sequence of
integers: $[t_1, t_2, \ldots, t_n]$ where each $t_i \in
\{0, 1, \ldots, V{-}1\}$ and $V$ is the vocabulary size.

% ----------------------------------------------------------------------------
\subsection{Embeddings: Words as Points in Space}
\label{sec:embeddings}
% ----------------------------------------------------------------------------

A token ID like 13389 is just a number---it tells us nothing about
what the token \emph{means}. The next step is to convert each token
into a \concept{high-dimensional vector}, called an
\concept{embedding}.

An embedding is a list of numbers (typically 768 to 4{,}096 of them)
that represents the token as a point in a high-dimensional space. The
critical property is that tokens with similar meanings end up at
nearby points:

\begin{figure}[htbp]
\centering
\begin{tikzpicture}[scale=1.0]
  % Axes
  \draw[-{Stealth}, thick, bitgray] (-0.3, 0) -- (8, 0);
  \draw[-{Stealth}, thick, bitgray] (0, -0.3) -- (0, 6);
  \node[font=\scriptsize\sffamily, bitgray, right] at (8, 0)
    {dimension 1};
  \node[font=\scriptsize\sffamily, bitgray, above] at (0, 6)
    {dimension 2};

  % Animal cluster
  \fill[bitblue!15, rounded corners=10pt]
    (0.5, 3.5) rectangle (3.5, 5.5);
  \node[font=\scriptsize\sffamily, bitblue] at (2, 5.8) {animals};
  \node[circle, fill=bitblue, inner sep=2pt, label={[font=\scriptsize]right:cat}]
    at (1.5, 4.8) {};
  \node[circle, fill=bitblue, inner sep=2pt, label={[font=\scriptsize]right:dog}]
    at (2.2, 4.3) {};
  \node[circle, fill=bitblue, inner sep=2pt, label={[font=\scriptsize]right:horse}]
    at (1.8, 3.9) {};

  % Vehicle cluster
  \fill[bitorange!15, rounded corners=10pt]
    (4.5, 3.5) rectangle (7.5, 5.5);
  \node[font=\scriptsize\sffamily, bitorange] at (6, 5.8) {vehicles};
  \node[circle, fill=bitorange, inner sep=2pt, label={[font=\scriptsize]right:car}]
    at (5.5, 4.8) {};
  \node[circle, fill=bitorange, inner sep=2pt, label={[font=\scriptsize]right:truck}]
    at (6.2, 4.3) {};
  \node[circle, fill=bitorange, inner sep=2pt, label={[font=\scriptsize]right:bus}]
    at (5.8, 3.9) {};

  % Royalty cluster
  \fill[bitpurple!15, rounded corners=10pt]
    (2.5, 0.5) rectangle (6.5, 2.5);
  \node[font=\scriptsize\sffamily, bitpurple] at (4.5, 2.8) {royalty};
  \node[circle, fill=bitpurple, inner sep=2pt, label={[font=\scriptsize]right:king}]
    at (3.5, 1.8) {};
  \node[circle, fill=bitpurple, inner sep=2pt, label={[font=\scriptsize]right:queen}]
    at (4.2, 1.3) {};
  \node[circle, fill=bitpurple, inner sep=2pt, label={[font=\scriptsize]right:prince}]
    at (5.5, 1.5) {};

  % Analogy arrow
  \draw[-{Stealth}, thick, dashed, bitgreen]
    (3.5, 1.8) -- node[font=\scriptsize\sffamily, above, sloped]
    {king $-$ man $+$ woman $\approx$ queen} (4.2, 1.3);

  % Note
  \node[font=\scriptsize\sffamily, bitgray, text width=4cm, align=center]
    at (8, 2.5) {(shown in 2D;\\actual embeddings\\use 768--4096\\dimensions)};
\end{tikzpicture}
\caption{Token embeddings place words in a high-dimensional space
  where distance reflects semantic similarity. Similar words cluster
  together, and vector arithmetic can capture analogies.
  This picture is simplified to 2D; real embeddings use hundreds to
  thousands of dimensions.}
\label{fig:embeddings}
\end{figure}

\begin{engineernote}
An embedding table is conceptually simple: it is a 2D array of shape
$V \times d$, where $V$ is the vocabulary size and $d$ is the
embedding dimension. Looking up a token's embedding is a single
table lookup---no computation, just a memory read of $d$ numbers.
In a model with $V = 50{,}000$ and $d = 4{,}096$, the embedding
table alone uses $50{,}000 \times 4{,}096 \times 2 = 400$\,MB
(in 16-bit precision).
\end{engineernote}

The embedding is not fixed forever. The embedding table is \emph{learned}
during training: the model adjusts the position of every token in this
high-dimensional space so that tokens appearing in similar contexts
end up nearby. This learned structure is called \concept{distributional
semantics}---the meaning of a word is defined by the company it keeps.

\subsubsection{Why High Dimensions Matter}

Two dimensions can separate cats from cars. But to simultaneously
encode that ``bank'' can mean a financial institution \emph{or} a
river bank, that ``run'' can be a verb \emph{or} a noun, and that
``Java'' can be a programming language \emph{or} an island---you need
many dimensions. Each dimension can encode a different aspect of
meaning, and a token's position in this space encodes all its
potential meanings simultaneously.

This is the first point of tension with bitwise approaches:
\emph{embeddings are inherently continuous}. The difference between
``happy'' and ``joyful'' is a small vector in this space; booleanising
it (rounding each dimension to 0 or~1) would destroy these fine
distinctions.

% ----------------------------------------------------------------------------
\subsection{The Transformer Block: The Engine of Modern AI}
\label{sec:transformer-block}
% ----------------------------------------------------------------------------

The transformer~\cite{vaswani2017attention} is the architecture
behind essentially all modern LLMs (GPT, Claude, Llama, etc.). A
transformer is a stack of identical \concept{transformer blocks}, each
containing two main sub-layers:

\begin{enumerate}[leftmargin=2em]
  \item \textbf{Self-attention}: lets each token look at every other
    token and decide which ones are relevant.
  \item \textbf{Feed-forward network (MLP)}: processes each token
    independently through a small neural network.
\end{enumerate}

Each sub-layer is wrapped in a \concept{residual connection} (add the
input back to the output) and \concept{layer normalisation} (rescale
the values to prevent them from growing too large or too small).

\begin{figure}[htbp]
\centering
\begin{tikzpicture}[
  block/.style={draw, rounded corners=3pt, minimum width=4cm,
    minimum height=0.8cm, font=\small\sffamily, align=center},
  arr/.style={-{Stealth[length=5pt]}, thick},
  plus/.style={circle, draw, minimum size=0.5cm, font=\small,
    fill=bitgreen!10, inner sep=0pt},
  skip/.style={thick, bitgray, dashed}
]
  % Input
  \node[font=\small\sffamily, bitgray] (in) at (0, 0) {token embeddings};

  % Self-attention
  \node[block, fill=bitblue!10, draw=bitblue] (attn)
    at (0, 1.5) {Self-Attention};
  \node[plus] (add1) at (0, 2.8) {$+$};
  \node[block, fill=bitgray!10, draw=bitgray] (ln1)
    at (0, 3.8) {Layer Normalisation};

  % MLP
  \node[block, fill=bitorange!10, draw=bitorange] (mlp)
    at (0, 5.0) {Feed-Forward Network (MLP)};
  \node[plus] (add2) at (0, 6.3) {$+$};
  \node[block, fill=bitgray!10, draw=bitgray] (ln2)
    at (0, 7.3) {Layer Normalisation};

  % Output
  \node[font=\small\sffamily, bitgray] (out) at (0, 8.5)
    {to next block};

  % Main flow
  \draw[arr] (in) -- (attn);
  \draw[arr] (attn) -- (add1);
  \draw[arr] (add1) -- (ln1);
  \draw[arr] (ln1) -- (mlp);
  \draw[arr] (mlp) -- (add2);
  \draw[arr] (add2) -- (ln2);
  \draw[arr] (ln2) -- (out);

  % Residual connections
  \draw[skip] (-2.8, 0.3) -- (-2.8, 2.8);
  \draw[skip, -{Stealth[length=4pt]}] (-2.8, 2.8) -- (add1);
  \draw[skip] (-2.8, 4.1) -- (-2.8, 6.3);
  \draw[skip, -{Stealth[length=4pt]}] (-2.8, 6.3) -- (add2);

  % Labels
  \node[font=\scriptsize\sffamily, bitgray, left] at (-3.2, 1.5)
    {residual};
  \node[font=\scriptsize\sffamily, bitgray, left] at (-3.2, 5.2)
    {residual};

  % Bracket for "one block"
  \draw[decorate, decoration={brace, amplitude=8pt, mirror},
    thick, bitpurple]
    (3, 0.7) -- (3, 8.0);
  \node[font=\small\sffamily, bitpurple, right, text width=3cm]
    at (3.5, 4.3)
    {one transformer\\block\\[4pt]\scriptsize (repeated\\32--80 times)};
\end{tikzpicture}
\caption{A single transformer block. The self-attention sub-layer lets
  tokens communicate with each other; the MLP sub-layer processes each
  token independently. Residual connections and layer normalisation
  stabilise training. A full LLM stacks 32--80 of these blocks.}
\label{fig:transformer-block}
\end{figure}

\begin{analogy}
Think of a transformer block as a two-step meeting process. In the
\emph{self-attention} step, every participant in the room looks
around, identifies who has relevant information, and takes notes.
In the \emph{MLP} step, each participant goes back to their desk
and privately thinks about what they learned. Then the whole process
repeats---another round of group discussion followed by private
reflection. After 32--80 rounds, each participant has a rich
understanding of the entire conversation.

The residual connection is like keeping your original notes alongside
the new ones---you never lose what you started with. Layer
normalisation is like agreeing on a common scale for everyone's
notes so that one loud participant does not dominate.
\end{analogy}

Let us examine each component in detail.

% ----------------------------------------------------------------------------
\subsection{Self-Attention: The Core Innovation}
\label{sec:self-attention}
% ----------------------------------------------------------------------------

Self-attention is the mechanism that lets each token in a sequence
attend to every other token. It answers the question: \emph{``For each
word I am processing, which other words in the sequence should
influence my understanding of it?''}

\subsubsection{The Intuition}

Consider the sentence: ``The animal didn't cross the street because
it was too wide.'' What does ``it'' refer to? A human instantly knows
``it'' means ``the street'' (streets are wide, not animals). But how
does a neural network figure this out?

Self-attention lets the token ``it'' \emph{look at every other token}
and assign each one a relevance score. It assigns high relevance to
``street'' and low relevance to ``animal,'' then uses these scores to
blend information from the relevant tokens into its own representation.

\subsubsection{The Mechanism: Queries, Keys, and Values}

Self-attention works through three learned projections. For each token,
the model computes:

\begin{itemize}[leftmargin=2em]
  \item A \concept{query} ($\bm{q}$): ``What am I looking for?''
  \item A \concept{key} ($\bm{k}$): ``What do I contain?''
  \item A \concept{value} ($\bm{v}$): ``What information do I
    provide if selected?''
\end{itemize}

\begin{figure}[htbp]
\centering
\begin{tikzpicture}[
  tok/.style={draw, rounded corners=2pt, minimum width=1.2cm,
    minimum height=0.6cm, font=\scriptsize\sffamily, fill=#1!10},
  mat/.style={draw, rounded corners=2pt, minimum width=1cm,
    minimum height=0.6cm, font=\scriptsize\sffamily, fill=#1!15},
  arr/.style={-{Stealth[length=4pt]}, thick},
  score/.style={font=\scriptsize\ttfamily}
]
  % Tokens
  \foreach \i/\w in {1/The, 2/cat, 3/sat, 4/on, 5/it} {
    \node[tok=bitgray] (t\i) at ({(\i-1)*2.5}, 0) {\w};
  }

  % Query from "it"
  \node[mat=bitblue, minimum width=1.5cm] (q) at (10, 1.5)
    {query: ``it''};
  \draw[arr, bitblue] (t5.north) -- (q.south);

  % Keys from all
  \foreach \i/\w in {1/The, 2/cat, 3/sat, 4/on, 5/it} {
    \node[mat=bitgreen] (k\i) at ({(\i-1)*2.5}, 2.5)
      {key: ``\w''};
    \draw[arr, bitgreen] (t\i.north) -- (k\i.south);
  }

  % Attention scores
  \node[font=\small\sffamily, bitgray] at (5, 4) {attention scores
    (query $\cdot$ key):};
  \foreach \i/\s/\h in {1/0.02/2, 2/0.71/35, 3/0.05/3, 4/0.03/2, 5/0.19/10} {
    \fill[bitblue!\h] ({(\i-1)*2.5-0.4}, 4.5) rectangle
      ({(\i-1)*2.5+0.4}, {4.5+\s*1.5});
    \node[score] at ({(\i-1)*2.5}, {4.5+\s*1.5+0.3}) {\s};
  }

  % Label
  \node[font=\scriptsize\sffamily, bitgray, text width=4cm, align=left]
    at (12.5, 3.5) {``it'' attends\\most strongly\\to ``cat''\\($\leftarrow$ correct\\referent)};
\end{tikzpicture}
\caption{Self-attention for the token ``it''. Its query is compared
  (via dot product) with every token's key to produce attention scores.
  The highest score (0.71) goes to ``cat''---the token ``it'' refers to.
  These scores then weight each token's value to produce the output.}
\label{fig:attention-scores}
\end{figure}

\subsubsection{The Mathematics}

Given a sequence of $n$ token embeddings packed into a matrix
$\bm{X} \in \mathbb{R}^{n \times d}$, self-attention computes:

\begin{align}
  \bm{Q} &= \bm{X}\bm{W}_Q, \quad
  \bm{K} = \bm{X}\bm{W}_K, \quad
  \bm{V} = \bm{X}\bm{W}_V
  \label{eq:qkv}\\[6pt]
  \text{Attention}(\bm{Q}, \bm{K}, \bm{V})
  &= \underbrace{\text{softmax}\!\left(
    \frac{\bm{Q}\bm{K}^\top}{\sqrt{d_k}}
  \right)}_{\text{attention weights}}
  \bm{V}
  \label{eq:attention}
\end{align}

where $\bm{W}_Q, \bm{W}_K, \bm{W}_V \in \mathbb{R}^{d \times d_k}$
are learned projection matrices, and $d_k$ is the dimension of each
head (typically $d / h$ where $h$ is the number of attention heads).

Let us unpack each step:

\begin{enumerate}[leftmargin=2em]
  \item \textbf{Project.} Multiply the input by three weight matrices
    to get queries, keys, and values. This is standard matrix
    multiplication---the same MAC operation we have been trying to
    eliminate.

  \item \textbf{Score.} Compute $\bm{Q}\bm{K}^\top$: the dot product
    of every query with every key. This produces an $n \times n$
    matrix of raw relevance scores. For a sequence of 4{,}096 tokens,
    this is a $4{,}096 \times 4{,}096$ matrix---about 16~million
    scores.

  \item \textbf{Scale.} Divide by $\sqrt{d_k}$ to prevent the dot
    products from growing too large (which would push softmax into
    regions where its gradient vanishes).

  \item \textbf{Softmax.} Apply the softmax function to each row:
    \begin{equation}
      \text{softmax}(z_i) = \frac{e^{z_i}}{\sum_j e^{z_j}}
      \label{eq:softmax}
    \end{equation}
    This converts raw scores into a probability distribution: all
    positive, summing to~1. The token with the highest relevance
    gets the most weight.

  \item \textbf{Aggregate.} Multiply the attention weights by the
    value matrix. Each token's output is a weighted average of all
    tokens' values, where the weights reflect relevance.
\end{enumerate}

\begin{keyinsight}
The softmax function is the critical bottleneck for binarisation.
It requires:
\begin{itemize}[leftmargin=1em]
  \item Exponentiation ($e^{z_i}$)---a transcendental function that
    maps integers to irrational numbers.
  \item Division---to normalise the probabilities.
  \item Continuous outputs---the weights 0.71 and 0.72 must be
    distinguishable, not rounded to the same bit.
\end{itemize}
Every other step (projection, scoring, aggregation) is matrix
multiplication, which \emph{can} be approximated with integer or
ternary arithmetic (as BitNet demonstrates). The softmax is where
floating-point computation concentrates.
\end{keyinsight}

\subsubsection{Multi-Head Attention}

In practice, the model does not compute one set of queries, keys, and
values. It computes $h$ sets in parallel (typically $h = 32$ to 128),
each with a smaller dimension $d_k = d/h$. These are called
\concept{attention heads}.

\begin{analogy}
Multi-head attention is like having multiple reading comprehension
specialists. One head might focus on syntactic relationships (subject
and verb agreement), another on semantic similarity (synonyms and
related concepts), a third on positional relationships (nearby words).
Each head independently decides what is relevant, and their
conclusions are concatenated and combined.
\end{analogy}

The output of multi-head attention is:
\begin{equation}
  \text{MultiHead}(\bm{X}) =
  [\text{head}_1; \text{head}_2; \ldots; \text{head}_h]\bm{W}_O
\end{equation}
where each $\text{head}_i = \text{Attention}(\bm{X}\bm{W}_Q^i,
\bm{X}\bm{W}_K^i, \bm{X}\bm{W}_V^i)$ and $\bm{W}_O \in
\mathbb{R}^{d \times d}$ is a learned output projection.

% ----------------------------------------------------------------------------
\subsection{The Feed-Forward Network (MLP Block)}
\label{sec:ffn}
% ----------------------------------------------------------------------------

After attention lets tokens communicate, each token is processed
\emph{independently} through a small two-layer neural network:

\begin{equation}
  \text{FFN}(\bm{x}) = \sigma(\bm{x}\bm{W}_1 + \bm{b}_1)\bm{W}_2 + \bm{b}_2
  \label{eq:ffn}
\end{equation}

where $\sigma$ is a nonlinear activation function (typically GELU or
SiLU), $\bm{W}_1 \in \mathbb{R}^{d \times d_{\text{ff}}}$ and
$\bm{W}_2 \in \mathbb{R}^{d_{\text{ff}} \times d}$.

The inner dimension $d_{\text{ff}}$ is typically $4 \times d$. For a
model with $d = 4{,}096$, this means $\bm{W}_1$ has $4{,}096 \times
16{,}384 = 67$~million parameters---\emph{per layer}.

\begin{engineernote}
The MLP block is where most of the model's parameters and computation
live. In a typical transformer, the MLP accounts for about two-thirds
of the total parameters and floating-point operations. This makes it
the highest-value target for binarisation: if you can make the MLP
bitwise, you have eliminated the majority of the floating-point work.
\end{engineernote}

\begin{analogy}
If self-attention is the ``group discussion'' where tokens share
information, the MLP is the ``private thinking'' where each token
processes what it learned. The attention layer moves information
\emph{between} tokens; the MLP transforms information \emph{within}
each token.
\end{analogy}

% ----------------------------------------------------------------------------
\subsection{Residual Connections and Layer Normalisation}
\label{sec:residual-layernorm}
% ----------------------------------------------------------------------------

Two seemingly minor components turn out to be essential for making
deep transformers trainable.

\subsubsection{Residual Connections}

A residual connection simply adds the input of a sub-layer back to
its output:
\begin{equation}
  \bm{y} = \bm{x} + \text{SubLayer}(\bm{x})
  \label{eq:residual}
\end{equation}

This means the sub-layer only needs to learn the \emph{difference}
from the input, not the entire transformation. Without residual
connections, gradients vanish in deep networks---a 96-layer
transformer would be untrainable.

\subsubsection{Layer Normalisation}

Layer normalisation~\cite{ba2016layernorm} rescales each token's
representation to have zero mean and unit variance:
\begin{equation}
  \text{LayerNorm}(\bm{x}) =
  \frac{\bm{x} - \mu}{\sqrt{\sigma^2 + \epsilon}} \odot \bm{\gamma} + \bm{\beta}
  \label{eq:layernorm}
\end{equation}
where $\mu$ and $\sigma^2$ are the mean and variance computed across
the embedding dimensions of a single token, $\bm{\gamma}$ and
$\bm{\beta}$ are learned scale and shift parameters, and $\epsilon$
is a small constant for numerical stability.

\begin{engineernote}
Layer normalisation is another floating-point bottleneck. It requires:
\begin{itemize}[leftmargin=1em]
  \item Computing mean and variance (requires division and square root).
  \item The division $(\bm{x} - \mu) / \sqrt{\sigma^2 + \epsilon}$
    produces continuous values even if $\bm{x}$ is integer.
  \item The learned parameters $\bm{\gamma}$ and $\bm{\beta}$ are
    continuous.
\end{itemize}
However, approximations exist: \concept{RMSNorm} (used in Llama and
many modern LLMs) drops the mean subtraction and simplifies to
$\bm{x} / \text{RMS}(\bm{x}) \cdot \bm{\gamma}$, which is slightly
simpler. Integer approximations of RMSNorm have been demonstrated
in quantisation research, though with some accuracy trade-off.
\end{engineernote}

% ----------------------------------------------------------------------------
\subsection{Autoregressive Generation: One Token at a Time}
\label{sec:autoregressive}
% ----------------------------------------------------------------------------

An LLM generates text through \concept{autoregressive generation}:
it predicts the next token, appends it to the sequence, and repeats.

\begin{figure}[htbp]
\centering
\begin{tikzpicture}[
  tok/.style={draw, rounded corners=2pt, minimum height=0.6cm,
    font=\scriptsize\sffamily, inner sep=3pt},
  pred/.style={tok, fill=bitgreen!20, draw=bitgreen},
  ctx/.style={tok, fill=bitblue!10, draw=bitblue},
  arr/.style={-{Stealth[length=4pt]}, thick, bitgray}
]
  % Step 1
  \node[font=\scriptsize\sffamily\bfseries, left] at (-0.5, 3)
    {Step 1:};
  \node[ctx] (s1a) at (1.5, 3) {The};
  \node[ctx] (s1b) at (3, 3) {cat};
  \node[pred] (s1c) at (4.5, 3) {sat};
  \draw[arr] (s1b) -- (s1c);

  % Step 2
  \node[font=\scriptsize\sffamily\bfseries, left] at (-0.5, 2)
    {Step 2:};
  \node[ctx] (s2a) at (1.5, 2) {The};
  \node[ctx] (s2b) at (3, 2) {cat};
  \node[ctx] (s2c) at (4.5, 2) {sat};
  \node[pred] (s2d) at (6, 2) {on};
  \draw[arr] (s2c) -- (s2d);

  % Step 3
  \node[font=\scriptsize\sffamily\bfseries, left] at (-0.5, 1)
    {Step 3:};
  \node[ctx] (s3a) at (1.5, 1) {The};
  \node[ctx] (s3b) at (3, 1) {cat};
  \node[ctx] (s3c) at (4.5, 1) {sat};
  \node[ctx] (s3d) at (6, 1) {on};
  \node[pred] (s3e) at (7.5, 1) {the};
  \draw[arr] (s3d) -- (s3e);

  % Step 4
  \node[font=\scriptsize\sffamily\bfseries, left] at (-0.5, 0)
    {Step 4:};
  \node[ctx] (s4a) at (1.5, 0) {The};
  \node[ctx] (s4b) at (3, 0) {cat};
  \node[ctx] (s4c) at (4.5, 0) {sat};
  \node[ctx] (s4d) at (6, 0) {on};
  \node[ctx] (s4e) at (7.5, 0) {the};
  \node[pred] (s4f) at (9, 0) {mat};
  \draw[arr] (s4e) -- (s4f);

  % Legend
  \node[ctx, right] at (10.5, 3) {context};
  \node[pred, right] at (10.5, 2.2) {predicted};
\end{tikzpicture}
\caption{Autoregressive generation. At each step, the model sees all
  previous tokens (the context) and predicts the next one. The
  predicted token is appended to the context, and the process repeats.
  Generating a 1{,}000-token response requires 1{,}000 forward passes.}
\label{fig:autoregressive}
\end{figure}

At each step, the model processes the entire context through all
transformer blocks and produces a probability distribution over the
vocabulary:

\begin{equation}
  P(t_{n+1} \mid t_1, \ldots, t_n) =
  \text{softmax}(\bm{h}_n \bm{W}_{\text{vocab}})
  \label{eq:next-token}
\end{equation}

where $\bm{h}_n$ is the hidden state of the last token after the
final transformer block, and $\bm{W}_{\text{vocab}} \in
\mathbb{R}^{d \times V}$ projects it into a score for each vocabulary
item.

\begin{keyinsight}
Autoregressive generation means the model's speed is determined by
\emph{latency per token}, not throughput. Even if you can process a
batch of tokens quickly (the ``prefill'' phase), generation is
inherently sequential: you cannot predict token~$n{+}1$ until you
have produced token~$n$. This is why LLM inference is
\emph{memory-bandwidth-bound}, not compute-bound---the bottleneck is
reading the model's weights from memory for each generated token, not
the arithmetic itself. This memory-bandwidth bottleneck is exactly
where bitwise models with smaller weights have the most to offer.
\end{keyinsight}

\subsubsection{The Causal Mask}

During generation, a token must only attend to tokens \emph{before}
it---it cannot look ahead at tokens that have not been generated yet.
This is enforced by a \concept{causal mask}: the attention scores for
future positions are set to $-\infty$ before the softmax, which drives
their attention weights to zero.

\begin{equation}
  \text{CausalAttention}(\bm{Q}, \bm{K}, \bm{V}) =
  \text{softmax}\!\left(
    \frac{\bm{Q}\bm{K}^\top}{\sqrt{d_k}} + \bm{M}
  \right)\bm{V}
\end{equation}

where $M_{ij} = 0$ if $i \geq j$ and $M_{ij} = -\infty$ otherwise.

% ----------------------------------------------------------------------------
\subsection{Putting It All Together: The Full LLM}
\label{sec:full-llm}
% ----------------------------------------------------------------------------

A complete LLM is assembled from the components we have described:

\begin{figure}[htbp]
\centering
\begin{tikzpicture}[
  comp/.style={draw, rounded corners=3pt, minimum width=5cm,
    minimum height=0.7cm, font=\small\sffamily, align=center},
  arr/.style={-{Stealth[length=5pt]}, thick},
  brace/.style={decorate, decoration={brace, amplitude=6pt, mirror}}
]
  % Stack
  \node[comp, fill=bitgray!8] (tok) at (0, 0)
    {Tokeniser: text $\to$ [token IDs]};
  \node[comp, fill=bitpurple!8] (emb) at (0, 1.2)
    {Embedding table: [IDs] $\to$ [vectors]};
  \node[comp, fill=bitblue!8] (b1) at (0, 2.4)
    {Transformer block 1};
  \node[comp, fill=bitblue!8] (b2) at (0, 3.4)
    {Transformer block 2};
  \node[font=\sffamily, bitgray] (dots) at (0, 4.2) {$\vdots$};
  \node[comp, fill=bitblue!8] (bn) at (0, 5.0)
    {Transformer block $N$};
  \node[comp, fill=bitgray!8] (ln) at (0, 6.2)
    {Final layer norm};
  \node[comp, fill=bitgreen!8] (head) at (0, 7.4)
    {Output head: vector $\to$ vocabulary scores};
  \node[comp, fill=bitorange!8] (sm) at (0, 8.6)
    {Softmax $\to$ next token probability};

  % Arrows
  \draw[arr] (tok) -- (emb);
  \draw[arr] (emb) -- (b1);
  \draw[arr] (b1) -- (b2);
  \draw[arr] (b2) -- (dots);
  \draw[arr] (dots) -- (bn);
  \draw[arr] (bn) -- (ln);
  \draw[arr] (ln) -- (head);
  \draw[arr] (head) -- (sm);

  % Annotations
  \draw[brace, thick, bitblue]
    (3.3, 2.1) -- (3.3, 5.3);
  \node[font=\scriptsize\sffamily, bitblue, right, text width=3.5cm]
    at (3.8, 3.7)
    {$N$ = 32 (7B model)\\$N$ = 80 (70B model)\\$N$ = 126 (405B model)};

  % Operation counts
  \node[font=\scriptsize\sffamily, bitgray, left, text width=3cm,
    align=right] at (-3.5, 1.2) {lookup only\\(no FP)};
  \node[font=\scriptsize\sffamily, bitgray, left, text width=3cm,
    align=right] at (-3.5, 3.7) {matrix multiply\\+ softmax\\+ layer norm};
  \node[font=\scriptsize\sffamily, bitgray, left, text width=3cm,
    align=right] at (-3.5, 8.6) {softmax\\(continuous)};
\end{tikzpicture}
\caption{The complete LLM architecture. Token IDs enter at the bottom,
  pass through $N$ transformer blocks, and produce a probability
  distribution over the next token at the top. The vast majority of
  computation happens in the transformer blocks (middle).}
\label{fig:full-llm}
\end{figure}

\Cref{tab:llm-sizes} shows the scale of modern LLMs.

\begin{table}[htbp]
\centering
\caption{Scale of representative LLMs. Parameter counts determine
  memory requirements; token counts determine the knowledge absorbed
  during training.}
\label{tab:llm-sizes}
\begin{tabular}{@{}lrrrrr@{}}
  \toprule
  \textbf{Model} & \textbf{Params} & \textbf{Layers} & \textbf{$d$}
    & \textbf{Heads} & \textbf{Training tokens} \\
  \midrule
  GPT-2       & 1.5B    & 48  & 1{,}600  & 25  & 40B \\
  Llama 2 7B  & 7B      & 32  & 4{,}096  & 32  & 2T \\
  Llama 2 70B & 70B     & 80  & 8{,}192  & 64  & 2T \\
  Llama 3 405B& 405B    & 126 & 16{,}384 & 128 & 15T \\
  \bottomrule
\end{tabular}
\end{table}

\begin{engineernote}
A 7-billion-parameter model stored in 16-bit floating point requires
$7 \times 10^9 \times 2 = 14$\,GB of memory---just for the weights.
Add activations, key-value cache, and optimizer state, and training
requires 4--8$\times$ more. This is why bitwise models are so
appealing: ternary weights ($\{-1, 0, +1\}$) need only 1.58~bits
per parameter, reducing the 7B model from 14\,GB to under 2\,GB.
\end{engineernote}

% ----------------------------------------------------------------------------
\subsection{Mixture of Experts: Scaling Without Proportional Cost}
\label{sec:moe}
% ----------------------------------------------------------------------------

As models grow larger, a key insight emerged: not every token needs
every parameter. \concept{Mixture of Experts} (MoE) replaces the
single MLP block in each transformer layer with multiple ``expert''
MLPs, and a small \concept{gating network} (also called a router)
decides which experts process each token.

\begin{figure}[htbp]
\centering
\begin{tikzpicture}[
  expert/.style={draw, rounded corners=3pt, minimum width=1.8cm,
    minimum height=1.2cm, font=\scriptsize\sffamily, align=center,
    fill=#1!10, draw=#1},
  gate/.style={draw, rounded corners=3pt, minimum width=2.5cm,
    minimum height=0.8cm, font=\small\sffamily, fill=bitgray!10},
  arr/.style={-{Stealth[length=4pt]}, thick},
  sel/.style={arr, bitgreen, line width=1.5pt},
  dim/.style={arr, bitgray!40, dashed}
]
  % Input
  \node[font=\small\sffamily] (in) at (5, -0.5) {token $\bm{x}$};

  % Gate
  \node[gate] (g) at (5, 1.2) {Router / Gate};
  \draw[arr] (in) -- (g);

  % Experts
  \node[expert=bitblue]   (e1) at (0,  3.5) {Expert 1\\(MLP)};
  \node[expert=bitgreen]  (e2) at (2.5, 3.5) {Expert 2\\(MLP)};
  \node[expert=bitorange] (e3) at (5,  3.5) {Expert 3\\(MLP)};
  \node[expert=bitpurple] (e4) at (7.5, 3.5) {Expert 4\\(MLP)};
  \node[expert=bitgray]   (e5) at (10, 3.5) {Expert 5\\(MLP)};

  % Gate to experts
  \draw[dim] (g.north) -- (e1.south);
  \draw[sel] (g.north) -- (e2.south);
  \draw[dim] (g.north) -- (e3.south);
  \draw[sel] (g.north) -- (e4.south);
  \draw[dim] (g.north) -- (e5.south);

  % Labels for selected
  \node[font=\scriptsize\sffamily, bitgreen, above right] at (e2.north)
    {selected};
  \node[font=\scriptsize\sffamily, bitgreen, above right] at (e4.north)
    {selected};

  % Output
  \node[font=\small\sffamily] (out) at (5, 5.5)
    {weighted sum of expert outputs};
  \draw[arr] (e2.north) -- (out.south west);
  \draw[arr] (e4.north) -- (out.south east);

  % Annotation
  \node[font=\scriptsize\sffamily, bitgray, text width=4cm, align=left]
    at (12, 2.5) {Top-$k$ routing:\\each token uses\\only 2 of 5 experts\\($\sim$40\% of params\\activated per token)};
\end{tikzpicture}
\caption{Mixture of Experts (MoE). A gating network routes each
  token to the top-$k$ experts (typically $k=2$). The model has many
  parameters (all experts combined) but each token only activates a
  fraction of them, keeping per-token computation manageable.}
\label{fig:moe}
\end{figure}

The gating function is typically a simple linear layer followed by
softmax:
\begin{equation}
  G(\bm{x}) = \text{softmax}(\bm{x}\bm{W}_g)
  \label{eq:gating}
\end{equation}
where $\bm{W}_g \in \mathbb{R}^{d \times E}$ and $E$ is the number
of experts. The top-$k$ values of $G(\bm{x})$ determine which experts
are activated and how their outputs are weighted.

\begin{keyinsight}
MoE decouples \emph{total model knowledge} from \emph{per-token
computation}. A model can have 400B total parameters but only activate
20B per token. This is relevant for bitwise approaches because:
\begin{itemize}[leftmargin=1em]
  \item The router/gating network is a small classification problem
    (``which expert should handle this token?'')---exactly the kind
    of task where Tsetlin Machines or Logic Gate Networks could excel.
  \item Each expert MLP is independent and could be individually
    compiled to logic gates or LUTs.
  \item The ``sparse activation'' principle aligns with bitwise
    networks' strength: small, cache-resident models that run fast.
\end{itemize}
\end{keyinsight}

\begin{analogy}
MoE is like a hospital with specialist departments. A patient (token)
does not see every doctor---a triage nurse (the router) assesses
symptoms and sends them to the relevant specialists. The hospital has
enormous total expertise (parameters), but each patient only uses a
small fraction. This is vastly more efficient than making every doctor
examine every patient.
\end{analogy}

% ----------------------------------------------------------------------------
\subsection{Where the Floating Point Lives}
\label{sec:where-fp-lives}
% ----------------------------------------------------------------------------

Now that we understand the full transformer architecture, we can
precisely identify where floating-point arithmetic is required and
where it might be eliminated. \Cref{tab:fp-map} provides this
breakdown.

\begin{table}[htbp]
\centering
\caption{Floating-point dependency map of a transformer. Recent
  research (I-BERT~\cite{kim2021ibert}, I-ViT~\cite{li2022ivit},
  I-LLM~\cite{hu2024illm}) has demonstrated integer-only
  approximations for components previously thought to require
  floating point.}
\label{tab:fp-map}
\begin{tabular}{@{}p{3.2cm}p{2.5cm}p{6.5cm}@{}}
  \toprule
  \textbf{Component} & \textbf{FP required?} & \textbf{Notes} \\
  \midrule
  Embedding lookup     & No  & Pure table lookup (integer index $\to$ vector). \\
  QKV projection       & Replaceable & Matrix multiply; BitNet uses ternary weights. \\
  $\bm{Q}\bm{K}^\top$ dot product & Replaceable & Integer dot product with ternary weights works. \\
  Softmax              & \textbf{Replaceable} & Integer approximations via polynomial~\cite{kim2021ibert} or bit-shift~\cite{li2022ivit} methods demonstrated. Softmax-free alternatives (linear attention, ReLU attention) also viable. See \cref{sec:softmax-problem}. \\
  Value aggregation    & Replaceable & Weighted sum; integer weights from approximate softmax suffice. \\
  Output projection    & Replaceable & Matrix multiply; can use ternary weights. \\
  Layer normalisation  & \textbf{Replaceable} & Integer RMSNorm via bit-shifting demonstrated in I-ViT and I-LLM~\cite{hu2024illm}. See \cref{sec:layernorm-problem}. \\
  MLP weights          & Replaceable & BitNet uses ternary $\{-1,0,+1\}$. \\
  MLP activation       & Replaceable & LUT-compiled or bit-shift approximation (ShiftGELU~\cite{li2022ivit}). \\
  MoE gating           & Replaceable & Small classification; softmax over few experts. \\
  \bottomrule
\end{tabular}
\end{table}

\begin{keyinsight}
A striking finding from recent research: \emph{every} component in
a transformer has a demonstrated integer-only or bitwise
approximation. The attention softmax can be replaced by polynomial
approximations~\cite{kim2021ibert}, bit-shift
operations~\cite{li2022ivit}, or eliminated entirely in favour of
softmax-free attention mechanisms. Layer normalisation can be
implemented with integer arithmetic and bit-shifts. The question is
no longer ``can we build an integer-only transformer?'' but
``how much accuracy do the combined approximations cost, and can
we close that gap?'' \Cref{sec:scaling-llms} explores this in
detail.
\end{keyinsight}
